% =============================================================================
%  Lab 0 - VHDL Fundamentals   *** EXAMPLE CHAPTER: Examples 1, 2 and 3 ***
%
%  Read this once, then write your own chapters in the same shape:
%    - every truth table has numbered rows
%    - every simplification names the law it uses
%    - code is included from code/ with \vhdlfile, never pasted
%    - every figure has a caption that says what it PROVES
%    - the Discussion describes a real mistake and what it taught
% =============================================================================

\chapter{Lab 0 --- VHDL Fundamentals}

\section{Objective}

Learn the structure of a VHDL design --- an \texttt{entity} that declares the
interface and an \texttt{architecture} that implements it --- by building and
verifying three small combinational circuits. Each circuit introduces one new way
of writing hardware:

\begin{table}[H]
\centering
\caption{The three circuits and what each one introduces.}
\label{tab:lab00-plan}
\begin{tabular}{cll}
\toprule
Ex. & circuit & new VHDL construct \\
\midrule
1 & \texttt{majority3} --- majority of three & concurrent Boolean assignment \\
2 & \texttt{mux2} --- 2-to-1 multiplexer & conditional assignment \texttt{when}/\texttt{else} \\
3 & \texttt{decoder2to4} --- 2-to-4 decoder with enable & vectors, \texttt{process}, \texttt{if}, \texttt{case} \\
\bottomrule
\end{tabular}
\end{table}

Every circuit is also verified with a self-checking testbench and a waveform, so
that correctness is proved by assertions rather than judged by eye.

% -----------------------------------------------------------------------------
\section{Specification}

\subsection{Example 1 --- \texttt{majority3}}

The output $Y$ is \texttt{1} when at least two of the three inputs $A$, $B$, $C$
are \texttt{1}. Three inputs give $2^3 = 8$ combinations, so the table is written
exhaustively (Table~\ref{tab:lab00-maj}).

\begin{table}[H]
\centering
\caption{Exhaustive truth table for \texttt{majority3}.}
\label{tab:lab00-maj}
\begin{tabular}{c|ccc|c|l}
\toprule
\# & $A$ & $B$ & $C$ & $Y$ & minterm \\
\midrule
0 & 0 & 0 & 0 & 0 & \\
1 & 0 & 0 & 1 & 0 & \\
2 & 0 & 1 & 0 & 0 & \\
3 & 0 & 1 & 1 & 1 & $m_3 = \overline{A}BC$ \\
4 & 1 & 0 & 0 & 0 & \\
5 & 1 & 0 & 1 & 1 & $m_5 = A\overline{B}C$ \\
6 & 1 & 1 & 0 & 1 & $m_6 = AB\overline{C}$ \\
7 & 1 & 1 & 1 & 1 & $m_7 = ABC$ \\
\bottomrule
\end{tabular}
\end{table}

\paragraph{Boolean derivation.}
\begin{align*}
\text{(1) canonical} \quad
  Y &= \textstyle\sum m(3,5,6,7) = \overline{A}BC + A\overline{B}C + AB\overline{C} + ABC \\[4pt]
\text{(2) simplified} \quad
  Y &= \overline{A}BC + A\overline{B}C + AB\overline{C} + ABC + ABC + ABC
     && \text{idempotent law: } X = X + X \\
  &= (\overline{A}+A)BC + A(\overline{B}+B)C + AB(\overline{C}+C)
     && \text{distributive law} \\
  &= BC + AC + AB
     && \text{complement law: } X + \overline{X} = 1 \\[4pt]
\text{(3) VHDL} \quad
  & \texttt{Y <= (A and B) or (A and C) or (B and C);}
\end{align*}

$Y$ does not appear on the right-hand side of its own equation: the circuit has no
memory, and its output depends only on the present inputs.

\subsection{Example 2 --- \texttt{mux2}}

The output $Y$ copies $D_0$ when the select input $S$ is \texttt{0}, and copies
$D_1$ when $S$ is \texttt{1}. Table~\ref{tab:lab00-muxE} lists all eight
combinations; Table~\ref{tab:lab00-muxC} states the same behaviour in compressed
form, naming the \emph{source} of the output instead of its value.

\begin{table}[H]
\centering
\begin{minipage}{0.52\linewidth}
\centering
\caption{Exhaustive table for \texttt{mux2}.}
\label{tab:lab00-muxE}
\begin{tabular}{c|ccc|c}
\toprule
\# & $S$ & $D_1$ & $D_0$ & $Y$ \\
\midrule
0 & 0 & 0 & 0 & 0 \\
1 & 0 & 0 & 1 & 1 \\
2 & 0 & 1 & 0 & 0 \\
3 & 0 & 1 & 1 & 1 \\
4 & 1 & 0 & 0 & 0 \\
5 & 1 & 0 & 1 & 0 \\
6 & 1 & 1 & 0 & 1 \\
7 & 1 & 1 & 1 & 1 \\
\bottomrule
\end{tabular}
\end{minipage}%
\begin{minipage}{0.44\linewidth}
\centering
\caption{Compressed table for \texttt{mux2}.}
\label{tab:lab00-muxC}
\begin{tabular}{c|c|c}
\toprule
\# & $S$ & $Y$ \\
\midrule
0 & 0 & $D_0$ \\
1 & 1 & $D_1$ \\
\bottomrule
\end{tabular}
\end{minipage}
\end{table}

\paragraph{Boolean derivation.}
\begin{align*}
\text{(1) canonical} \quad
  Y &= \textstyle\sum m(1,3,6,7)
     = \overline{S}\,\overline{D_1}D_0 + \overline{S}D_1D_0 + SD_1\overline{D_0} + SD_1D_0 \\[4pt]
\text{(2) simplified} \quad
  Y &= \overline{S}D_0(\overline{D_1} + D_1) + SD_1(\overline{D_0} + D_0)
     && \text{distributive law} \\
  &= \overline{S}D_0 + SD_1
     && \text{complement law} \\[4pt]
\text{(3) VHDL} \quad
  & \texttt{Y <= D0 when S = '0' else D1;}
\end{align*}

Read aloud, the simplified equation is the specification itself: $D_0$ passes when
$S$ is low, $D_1$ passes when $S$ is high.

\subsection{Example 3 --- \texttt{decoder2to4}}

When the enable $E$ is \texttt{1}, exactly one of the four outputs $D_3 \dots D_0$
is \texttt{1}: the one whose index equals the binary address $A_1A_0$. When $E$ is
\texttt{0}, all outputs are \texttt{0}. Row~0 of Table~\ref{tab:lab00-dec} uses
don't-cares (--) to state that the address has no effect while the decoder is
disabled.

\begin{table}[H]
\centering
\caption{Truth table for \texttt{decoder2to4}. Rows 1--4 are one-hot.}
\label{tab:lab00-dec}
\begin{tabular}{c|ccc|cccc|l}
\toprule
\# & $E$ & $A_1$ & $A_0$ & $D_3$ & $D_2$ & $D_1$ & $D_0$ & note \\
\midrule
0 & 0 & -- & -- & 0 & 0 & 0 & 0 & disabled \\
1 & 1 & 0 & 0 & 0 & 0 & 0 & 1 & \\
2 & 1 & 0 & 1 & 0 & 0 & 1 & 0 & \\
3 & 1 & 1 & 0 & 0 & 1 & 0 & 0 & \\
4 & 1 & 1 & 1 & 1 & 0 & 0 & 0 & \\
\bottomrule
\end{tabular}
\end{table}

\paragraph{Boolean equations.} There is one equation per output. Each output is a
single product term --- one minterm of the address, gated by the enable --- so
nothing simplifies further:
\begin{align*}
D_0 &= E\,\overline{A_1}\,\overline{A_0} &
D_1 &= E\,\overline{A_1}\,A_0 &
D_2 &= E\,A_1\,\overline{A_0} &
D_3 &= E\,A_1\,A_0
\end{align*}
or, compactly, $D_k = E \cdot m_k(A_1, A_0)$.

% -----------------------------------------------------------------------------
\section{Design}

The same circuit can be described in three views. Each view is useful for a
different purpose, and part of designing is choosing which one to write.

\begin{table}[H]
\centering
\caption{Three views of one design.}
\label{tab:lab00-views}
\small
\begin{tabular}{l>{\raggedright\arraybackslash}p{6.2cm}>{\raggedright\arraybackslash}p{4.8cm}}
\toprule
view & what you write & best when \\
\midrule
behavioural & \emph{what} the output is (equation, condition, case) & the specification is a rule \\
structural  & \emph{how} it is wired: one signal or component per block & you reuse parts or care about the wiring \\
schematic   & a drawing of gates and wires & explaining and checking the structure \\
\bottomrule
\end{tabular}
\end{table}

Quartus accepts the first two as VHDL and produces the third automatically in the
RTL Viewer. Comparing your own schematic with the RTL Viewer is the fastest way to
check that the tool understood the design the way you intended.

\subsection{Example 1 --- three views of \texttt{majority3}}

\begin{itemize}[leftmargin=1.5em]
\item \textbf{Behavioural:} the simplified equation written as one concurrent
assignment, \texttt{Y <= (A and B) or (A and C) or (B and C);}
\item \textbf{Structural (gate level):} one internal signal per gate --- three AND
outputs \texttt{AB\_s}, \texttt{AC\_s}, \texttt{BC\_s}, then one OR.
\item \textbf{Schematic:} Figure~\ref{fig:lab00-maj-sch}. Each product term of the
equation is one AND gate; the sum is the OR gate.
\end{itemize}

\begin{figure}[H]
\centering
% Schematic view of majority3: Y = AB + AC + BC
% Vertical lines are input buses; a dot marks a connection, a crossing without
% a dot is not connected.
\begin{circuitikz}
  \coordinate (A) at (-3.0, 2.3);
  \coordinate (B) at (-2.3, 2.3);
  \coordinate (C) at (-1.6, 2.3);
  \node[and port] (g1) at (0, 1.5) {};
  \node[and port] (g2) at (0, 0)   {};
  \node[and port] (g3) at (0,-1.5) {};
  \node[or port, number inputs=3] (g4) at (3.2, 0) {};
  \draw (A) node[above] {$A$} -- (A |- 0,-2.2);
  \draw (B) node[above] {$B$} -- (B |- 0,-2.2);
  \draw (C) node[above] {$C$} -- (C |- 0,-2.2);
  \draw (g1.in 1) -- (g1.in 1 -| A) node[circ] {};
  \draw (g1.in 2) -- (g1.in 2 -| B) node[circ] {};
  \draw (g2.in 1) -- (g2.in 1 -| A) node[circ] {};
  \draw (g2.in 2) -- (g2.in 2 -| C) node[circ] {};
  \draw (g3.in 1) -- (g3.in 1 -| B) node[circ] {};
  \draw (g3.in 2) -- (g3.in 2 -| C) node[circ] {};
  \draw (g1.out) node[above right] {\small $AB$} -| (g4.in 1);
  \draw (g2.out) node[above right] {\small $AC$} -- (g4.in 2);
  \draw (g3.out) node[below right] {\small $BC$} -| (g4.in 3);
  \draw (g4.out) -- ++(0.5,0) node[right] {$Y$};
\end{circuitikz}

\caption{Schematic view of \texttt{majority3}: three 2-input AND gates, one per
product term of $Y = AB + AC + BC$, feeding a 3-input OR gate. A dot marks a
connection; lines that cross without a dot are not connected.}
\label{fig:lab00-maj-sch}
\end{figure}

\textbf{Chosen view: behavioural.} The simplified equation maps directly onto one
line of VHDL. The structural version adds three internal signals without changing
the synthesised hardware, so it is written only to show the correspondence between
the equation, the code and the schematic.

\subsection{Example 2 --- \texttt{mux2}}

The specification is a \emph{choice} (``if $S$ is 0, pass $D_0$''), so the
behavioural \texttt{when}/\texttt{else} form states it most directly. The equation
$Y = \overline{S}D_0 + SD_1$ gives the schematic in Figure~\ref{fig:lab00-mux-sch}:
an inverter produces $\overline{S}$, and each AND gate lets one data input through.

\begin{figure}[H]
\centering
% Schematic view of mux2: Y = S'D0 + S D1
\begin{circuitikz}
  \coordinate (D0) at (-3.0, 2.0);
  \coordinate (D1) at (-2.3, 2.0);
  \coordinate (S)  at (-1.6, 2.0);
  \node[not port, scale=0.7] (n1) at (-0.9, -1.9) {};
  \node[and port] (g1) at (0.9, 1.0)  {};
  \node[and port] (g2) at (0.9, -1.0) {};
  \node[or port]  (g3) at (3.4, 0)    {};
  \draw (D0) node[above] {$D_0$} -- (D0 |- 0,-2.4);
  \draw (D1) node[above] {$D_1$} -- (D1 |- 0,-2.4);
  \draw (S)  node[above] {$S$}   -- (S |- 0,-2.4);
  \draw (n1.in) -- (n1.in -| S) node[circ] {};
  \draw (g1.in 1) -- (g1.in 1 -| D0) node[circ] {};
  \draw (g1.in 2) -- (g1.in 2 -| n1.out) -- (n1.out);
  \draw (g2.in 1) -- (g2.in 1 -| D1) node[circ] {};
  \draw (g2.in 2) -- (g2.in 2 -| S) node[circ] {};
  \node[right] at (n1.out) {\small $\overline{S}$};
  \draw (g1.out) node[above right] {\small $\overline{S}D_0$} -| (g3.in 1);
  \draw (g2.out) node[below right] {\small $SD_1$} -| (g3.in 2);
  \draw (g3.out) -- ++(0.5,0) node[right] {$Y$};
\end{circuitikz}

\caption{Schematic view of \texttt{mux2}. Exactly one AND gate is enabled for any
value of $S$, so exactly one data input reaches $Y$.}
\label{fig:lab00-mux-sch}
\end{figure}

\subsection{Example 3 --- four ways to write \texttt{decoder2to4}}

The decoder is the first circuit with a vector output, and it can be written in at
least four equivalent forms:

\begin{table}[H]
\centering
\caption{Four VHDL forms of the same decoder.}
\label{tab:lab00-decforms}
\small
\begin{tabular}{>{\raggedright\arraybackslash}p{3.6cm}>{\raggedright\arraybackslash}p{6.4cm}>{\raggedright\arraybackslash}p{4.4cm}}
\toprule
form & VHDL construct & enable handled by \\
\midrule
(a) Boolean equations & four concurrent assignments \texttt{D(k) <= E and \dots} & the \texttt{E} term in each product \\
(b) selected assignment & \texttt{with A select D <= \dots} & must be added explicitly \\
(c) process with \texttt{if}/\texttt{case} & \texttt{if E = '0' \dots else case A is \dots} & the \texttt{if}, before the \texttt{case} \\
(d) process with default & \texttt{D <= "0000"; if E = '1' then case \dots} & the default assignment \\
\bottomrule
\end{tabular}
\end{table}

Form (a) is the schematic written as text (Figure~\ref{fig:lab00-dec-sch}): two
inverters produce $\overline{A_1}$ and $\overline{A_0}$, and each 3-input AND gate
computes one minterm gated by $E$.

\begin{figure}[htbp]
\centering
% Schematic view of decoder2to4: D(k) = E . m_k(A1, A0)
\begin{circuitikz}
  \coordinate (E)   at (-4.2, 3.4);
  \coordinate (A1)  at (-3.5, 3.4);
  \coordinate (nA1) at (-2.6, 3.4);
  \coordinate (A0)  at (-1.9, 3.4);
  \coordinate (nA0) at (-1.0, 3.4);
  \node[and port, number inputs=3] (d0) at (1.2,  2.1) {};
  \node[and port, number inputs=3] (d1) at (1.2,  0.5) {};
  \node[and port, number inputs=3] (d2) at (1.2, -1.1) {};
  \node[and port, number inputs=3] (d3) at (1.2, -2.7) {};
  \draw (E)  node[above] {$E$}   -- (E  |- 0,-3.4);
  \draw (A1) node[above] {$A_1$} -- (A1 |- 0,-3.4);
  \draw (A0) node[above] {$A_0$} -- (A0 |- 0,-3.4);
  \node[not port, scale=0.55, rotate=-90] (i1) at ($(nA1)+(0,-0.35)$) {};
  \node[not port, scale=0.55, rotate=-90] (i0) at ($(nA0)+(0,-0.35)$) {};
  \draw (A1 |- 0,3.05) node[circ] {} -- ++(0,0.3) -| (i1.in);
  \draw (A0 |- 0,3.05) node[circ] {} -- ++(0,0.3) -| (i0.in);
  \draw (i1.out) -- (nA1 |- 0,-3.4);
  \draw (i0.out) -- (nA0 |- 0,-3.4);
  \node[above] at ($(nA1)+(0,0.35)$) {$\overline{A_1}$};
  \node[above] at ($(nA0)+(0,0.35)$) {$\overline{A_0}$};
  \foreach \g in {d0,d1,d2,d3} { \draw (\g.in 1) -- (\g.in 1 -| E) node[circ] {}; }
  \draw (d0.in 2) -- (d0.in 2 -| nA1) node[circ] {};  \draw (d0.in 3) -- (d0.in 3 -| nA0) node[circ] {};
  \draw (d1.in 2) -- (d1.in 2 -| nA1) node[circ] {};  \draw (d1.in 3) -- (d1.in 3 -| A0)  node[circ] {};
  \draw (d2.in 2) -- (d2.in 2 -| A1)  node[circ] {};  \draw (d2.in 3) -- (d2.in 3 -| nA0) node[circ] {};
  \draw (d3.in 2) -- (d3.in 2 -| A1)  node[circ] {};  \draw (d3.in 3) -- (d3.in 3 -| A0)  node[circ] {};
  \draw (d0.out) -- ++(0.5,0) node[right] {$D_0$};
  \draw (d1.out) -- ++(0.5,0) node[right] {$D_1$};
  \draw (d2.out) -- ++(0.5,0) node[right] {$D_2$};
  \draw (d3.out) -- ++(0.5,0) node[right] {$D_3$};
\end{circuitikz}

\caption{Schematic view of \texttt{decoder2to4}. $E$ enters every AND gate, so when
$E = 0$ all four outputs are forced to 0 regardless of the address.}
\label{fig:lab00-dec-sch}
\end{figure}

\textbf{Chosen form: (c).} The specification has a priority --- \emph{first} check
the enable, \emph{then} decode the address --- and the \texttt{if} followed by a
\texttt{case} states that priority in the same order. Form (b) is shorter but has no
natural place for the enable; Section~\ref{sec:lab00-discussion} describes what
happened when the first version used it.

A \texttt{process} does not by itself mean sequential hardware. This process
describes combinational logic because it has no clock and \texttt{D} receives a
value on every path through it.

% -----------------------------------------------------------------------------
\section{Implementation}

\vhdlfile{code/lab00/majority3.vhd}{\texttt{majority3} --- majority of three (behavioural)}

\vhdlfile{code/lab00/mux2.vhd}{\texttt{mux2} --- 2-to-1 multiplexer (behavioural)}

\vhdlfile{code/lab00/decoder2to4.vhd}{\texttt{decoder2to4} --- 2-to-4 decoder with enable, corrected version (form c)}

% -----------------------------------------------------------------------------
\section{Verification}

All three testbenches follow the same pattern: the testbench entity has no ports,
it instantiates the design under test with named association, applies one vector
every 10\,ns, \emph{checks} the output after each wait, and prints a single
\texttt{PASS} or \texttt{FAILED} verdict at the end.

\paragraph{Coverage argument.}
\begin{itemize}[leftmargin=1.5em]
\item \textbf{\texttt{majority3}:} all 8 rows of Table~\ref{tab:lab00-maj} are
applied, stored in the testbench as a constant table. Exhaustive testing leaves no
untested input combination.
\item \textbf{\texttt{mux2}:} 4 vectors, two for each value of $S$. In every vector
$D_0 \neq D_1$, so the output is only correct if the multiplexer selected the right
input; a circuit that ignored $S$ or ORed its inputs would fail.
\item \textbf{\texttt{decoder2to4}:} all 5 rows of Table~\ref{tab:lab00-dec}. Row~0
($E = 0$) is applied first, because the enable is part of the specification and a
testbench that only drives $E = 1$ cannot detect a decoder that ignores it.
\end{itemize}

\vhdlfile{code/lab00/majority3_tb.vhd}{\texttt{majority3\_tb} --- exhaustive, table-driven}

\vhdlfile{code/lab00/mux2_tb.vhd}{\texttt{mux2\_tb} --- four vectors with a \texttt{check} procedure}

\vhdlfile{code/lab00/decoder2to4_tb.vhd}{\texttt{decoder2to4\_tb} --- all five table rows, disabled row first}

% -----------------------------------------------------------------------------
\section{Results}

\subsection{Compilation}

All designs compiled with 0 errors on device \texttt{5CSEMA4U23C6}
(Table~\ref{tab:lab00-quartus}).

\begin{table}[H]
\centering
\caption{Quartus results. Pin counts equal inputs plus outputs; no design contains
a register, as expected for combinational logic.}
\label{tab:lab00-quartus}
\begin{tabular}{lcccc}
\toprule
& \texttt{majority3} & \texttt{mux2} & \texttt{decoder2to4} v1 & \texttt{decoder2to4} \\
\midrule
pins (in / out) & 4 (3 / 1) & 4 (3 / 1) & 7 (3 / 4) & 7 (3 / 4) \\
logic (ALMs)    & 1 & 1 & 3 & 3 \\
registers       & 0 & 0 & 0 & 0 \\
design warning  & none & none & \textbf{15610} & none \\
\bottomrule
\end{tabular}
\end{table}

Three messages appear for every design and are expected, not errors:
\textit{Critical Warning 169085} (no pin locations --- there is no board this
semester), \textit{Critical Warning 332012} (no timing constraints file) and
\textit{Warning 332068} (no clocks). A combinational circuit has no clock, so the
Timing Analyzer reports no $F_{max}$.

The first version of the decoder produced one additional warning, which is the bug
discussed in Section~\ref{sec:lab00-discussion}:
\begin{lstlisting}[language={},numbers=none,caption={}]
Warning (21074): Design contains 1 input pin(s) that do not drive logic
    Warning (15610): No output dependent on input pin "E"
\end{lstlisting}

\subsection{Simulation}

Each vector occupies one 10\,ns window, so row $i$ of a table is visible between
$10i$ and $10i+10$\,ns. Each waveform figure shows the signal rows of the Questa
screenshot directly above its time axis; the empty area between them is not shown.

\wavefig{figures/majority3_waveform.png}{408bp}{447bp}{\texttt{majority3\_tb} waveform. $Y$ is
high only in the windows 30--40\,ns and 50--80\,ns, which are rows 3, 5, 6 and 7 of
Table~\ref{tab:lab00-maj}: the output is 1 exactly when two or more inputs are 1.
The run ends at 80\,ns, confirming all 8 rows were applied.}{fig:lab00-maj-wave}

\wavefig{figures/2t1_multiplexer_waveform.png}{407bp}{445bp}{\texttt{mux2\_tb} waveform. In
0--10\,ns and 20--30\,ns ($S=0$) $Y$ equals $D_0$ (0, then 1); in 10--20\,ns and
30--40\,ns ($S=1$) $Y$ equals $D_1$ (1, then 0). Because $D_0 \neq D_1$ in every
window, $Y$ changing with $S$ proves the select input chooses the
source.}{fig:lab00-mux-wave}

\wavefig{figures/decoder_waveform.png}{408bp}{446bp}{\texttt{decoder2to4\_tb} waveform of the
\emph{first} version (\texttt{D} shown in hexadecimal). In 10--50\,ns the outputs
step through 1, 2, 4, 8 --- one-hot, rows 1--4 correct. In 0--10\,ns, however,
$E = 0$ and $D = 1$ (\texttt{0001}): row~0 of Table~\ref{tab:lab00-dec} is violated,
because this version ignores the enable.}{fig:lab00-dec-wave}

\begin{table}[H]
\centering
\caption{Simulation transcripts (Questa).}
\label{tab:lab00-transcripts}
\begin{tabular}{lll}
\toprule
testbench & design & transcript verdict \\
\midrule
\texttt{majority3\_tb}   & \texttt{majority3}      & \texttt{PASS - majority3: all 8 rows correct} (80\,ns) \\
\texttt{mux2\_tb}        & \texttt{mux2}           & \texttt{PASS - mux2: all 4 vectors correct} (40\,ns) \\
\texttt{decoder2to4\_tb} & first version  & \texttt{FAIL: row 0: E=0 must disable every output} (10\,ns) \\
\texttt{decoder2to4\_tb} & corrected      & \texttt{PASS - decoder2to4: all 5 rows correct} (50\,ns) \\
\bottomrule
\end{tabular}
\end{table}

The corrected decoder differs from Figure~\ref{fig:lab00-dec-wave} only in the
0--10\,ns window, where $D$ is now \texttt{0000}.

% -----------------------------------------------------------------------------
\section{Discussion}
\label{sec:lab00-discussion}

\paragraph{The decoder ignored its enable.} The first version of the decoder used a
selected signal assignment:

\vhdlfile{code/lab00/decoder2to4_v1.vhd}{First version of \texttt{decoder2to4} --- \texttt{E} is declared but never used}

It lists all four addresses correctly, and rows 1--4 of the waveform were right.
But \texttt{E} appears only in the entity: the architecture never reads it, so the
circuit decodes the address whether or not it is enabled. The final
\texttt{"0000" when others} line looks as if it handles the disabled case, but it
does not --- it only covers address values such as \texttt{'U'} or \texttt{'X'},
which exist in simulation and not in hardware.

The bug was visible in three independent places:
\begin{enumerate}[leftmargin=1.8em]
\item \textbf{Quartus messages:} \textit{Warning 15610: No output dependent on input
pin "E"} --- the tool reported that an input drives nothing.
\item \textbf{Waveform:} $D = \texttt{0001}$ while $E = 0$ in the first 10\,ns
(Figure~\ref{fig:lab00-dec-wave}).
\item \textbf{Assertion:} the row-0 check failed at 10\,ns.
\end{enumerate}

It was fixed by rewriting the architecture as form~(c): test \texttt{E} first, then
decode the address with a \texttt{case}. After the change the Quartus warning
disappeared and the testbench reported \texttt{PASS}.

\paragraph{What this taught.} A testbench can only find a bug in a row it actually
applies. If the testbench had started with $E = 1$, all of its checks would have
passed and the broken decoder would have looked verified. Every row of the truth
table --- including the ``boring'' disabled row --- must become a vector. And the
warnings in the Messages window are part of the evidence: this one described the
bug exactly, before any simulation was run.
